%%% SECTION 6: ADAPTED 21 CFR PART 312 FOR PHYSICAL AI IND APPLICATIONS %%%

[This section should be approximately 10-12 pages. The major theme is that
Physical AI Adaption: 21 CFR Part 312 \cite{cfr312-adapt} provides a
comprehensive framework for Investigational New Drug applications in the
context of Physical AI oncology trials, building on the existing regulation
\cite{cfr-part-312}. This is the most extensive of the three regulatory
adaptations, spanning ten subparts including a new Subpart J dedicated entirely
to Physical AI systems. The adaptation demonstrates that Physical AI can be
integrated into the IND process without requiring entirely new legislation.]

[PRIMARY SOURCE: national-platform/21cfr312\_adapt/ directory containing five
chunked files: 01\_preamble\_scope\_definitions.tex (lines 1-427),
02\_ind\_content\_phases.tex (lines 428-701),
03\_protocol\_amendments\_reporting.tex (lines 702-1029),
04\_annual\_reports\_withdrawal.tex (lines 1030-1492), and
05\_clinical\_holds\_appendices\_closing.tex (lines 1493-2275). All five files
must be used together to build this section.]

\subsection{Scope and Physical AI Definitions}
\label{subsec:cfr312-scope}

[Build from national-platform/21cfr312\_adapt/01\_preamble\_scope\_definitions.tex,
Subpart A (General Provisions, \S 312.1-312.10). Cover scope (\S 312.1),
applicability (\S 312.2), and the 22 Physical AI definitions (\S 312.3).
These definitions are the most comprehensive set across all three adapted
standards. Include definitions for Physical AI labeling (\S 312.6), Physical AI
promotion restrictions (\S 312.7), Physical AI charging provisions (\S 312.8),
and waiver procedures (\S 312.10).]

\subsection{IND Requirements and Physical AI Content}
\label{subsec:cfr312-ind}

[Build from national-platform/21cfr312\_adapt/02\_ind\_content\_phases.tex,
Subpart B (IND Application). Cover:]

\subsubsection{IND Requirement for Physical AI Trials}
[Detail \S 312.20 on when an IND is required for Physical AI trials and how
the requirement extends to trials where Physical AI systems are integral to
drug administration or monitoring.]

\subsubsection{Phases of Investigation}
[Detail \S 312.21 including the adapted clinical trial phases with special
attention to the Phase 0 simulation validation stage. This is a significant
innovation: Physical AI trials include a mandatory simulation phase before any
human subjects are enrolled, using tools like those in the main repository
\cite{main-repo}.]

[MAJOR THEME: The Phase 0 simulation validation requirement is directly
supported by the evidence from both A Cancer Patient's Journey, Physical AI
\cite{patient-journey-paper} (single-patient simulation) and the 24-hour
simulation from Clinical Trial Site Complete Documentation \cite{site-docs}
(168-patient simulation). These completed simulations demonstrate that
Physical AI systems can be thoroughly validated before real patient contact.]

\subsubsection{IND Content and Format}
[Detail \S 312.23 including the Physical AI System Description requirement
for IND submissions. Cover what information about Physical AI systems must
be included in the IND application.]

\subsubsection{Protocol Amendments}
[Detail \S 312.30 including Physical AI amendment triggers: when changes to
Physical AI systems require protocol amendments, notification to FDA, and
IRB approval.]

\subsection{Safety Reporting for Physical AI}
\label{subsec:cfr312-safety}

[Build from national-platform/21cfr312\_adapt/03\_protocol\_amendments\_reporting.tex.
Cover:]

\subsubsection{Information Amendments}
[Detail \S 312.31 on information amendments specific to Physical AI system
updates and modifications.]

\subsubsection{IND Safety Reporting}
[Detail \S 312.32 including Physical AI adverse event reporting and
cybersecurity incident reporting. This subsection establishes reporting
timelines and severity classifications for robot malfunctions, AI decision
errors, sensor failures, and cybersecurity breaches that may affect patient
safety.]

\subsubsection{Annual Reports}
[Detail \S 312.33 including the Physical AI performance summary required in
annual reports, covering system uptime, adverse events, software updates, and
robot maintenance records.]

\subsubsection{IND Withdrawal}
[Detail \S 312.38 including Physical AI system decommissioning procedures when
an IND is withdrawn.]

\subsection{Administrative Actions for Physical AI Trials}
\label{subsec:cfr312-admin}

[Build from national-platform/21cfr312\_adapt/03\_protocol\_amendments\_reporting.tex
and 04\_annual\_reports\_withdrawal.tex, Subpart C. Cover general requirements
(\S 312.40), termination including Physical AI grounds (\S 312.44), inactive
status and Physical AI dormancy/reactivation (\S 312.45), meetings including
Physical AI system review (\S 312.47), and dispute resolution for Physical AI
technical disputes (\S 312.48).]

\subsection{Sponsor and Investigator Responsibilities}
\label{subsec:cfr312-responsibilities}

[Build from national-platform/21cfr312\_adapt/04\_annual\_reports\_withdrawal.tex,
Subpart D (\S 312.50-312.70). Cover all Physical AI adaptations to sponsor and
investigator responsibilities. Cross-reference with the ICH E6(R3) adaptation
(Section~\ref{sec:ich-e6r3}) for complementary requirements. Emphasize that
both sponsors and investigators must have Physical AI competencies and that
the Adaption: ICH Harmonised Guideline \cite{ich-adapt} provides detailed
training requirements.]

\subsection{Life-Threatening Illness Provisions}
\label{subsec:cfr312-life-threatening}

[Build from national-platform/21cfr312\_adapt/05\_clinical\_holds\_appendices\_closing.tex,
Subpart E (\S 312.80-312.88). Cover how expanded access and accelerated
approval provisions apply to Physical AI oncology treatments for
life-threatening cancers. Note that oncology is a primary use case for these
provisions, as many cancer patients face life-threatening conditions where
Physical AI may offer faster or better treatment options.]

\subsection{Expanded Access for Physical AI}
\label{subsec:cfr312-expanded-access}

[Build from national-platform/21cfr312\_adapt/05\_clinical\_holds\_appendices\_closing.tex,
Subpart I (\S 312.300-312.320). Cover expanded access provisions adapted for
Physical AI, including how Physical AI systems can be made available outside
of clinical trials when sufficient evidence supports their safety and
effectiveness.]

\subsection{Subpart J: Physical AI Systems}
\label{subsec:cfr312-subpart-j}

[Build from national-platform/21cfr312\_adapt/05\_clinical\_holds\_appendices\_closing.tex,
Subpart J (\S 312.400-312.405). This is entirely new and dedicated to Physical
AI systems. Cover all provisions in this new subpart, which establishes
Physical AI-specific regulatory requirements that do not fit within the
existing subpart structure.]

[MAJOR THEME: The creation of a dedicated Subpart J for Physical AI demonstrates
that while the existing CFR Part 312 framework can accommodate Physical AI
through adaptation, certain aspects of Physical AI are sufficiently novel to
warrant their own regulatory provisions. This balanced approach of adaptation
plus targeted new provisions represents the most practical path to regulatory
readiness.]
